\documentclass[10pt,a4paper]{amsart}
\usepackage[margin=19mm]{geometry}
\usepackage{amsmath,amssymb,mathtools}
\usepackage[scr=rsfs,cal=euler]{mathalfa}
\usepackage{xfrac}
\usepackage[unicode,hidelinks,bookmarks=false]{hyperref}
\newcommand{\ie}{i.e.\ }
\newcommand{\cf}{cf.\ }
\newcommand{\resp}{respectively\ }
\newcommand{\Subsection}{\subsection}
\providecommand{\nobreakdash}{\nobreak\hbox{-}}
\setlength{\emergencystretch}{2em}
\multlinegap0pt
\usepackage{amssymb}
\usepackage{bm}
\usepackage{enumerate}
\usepackage{mathtools}
\usepackage{tikz}
\newtheorem{lemma}{Lemma}
\newtheorem{definition}{Definition}
\newtheorem{prop}{Proposition}
\newcommand{\C}{\mathbb{C}}           % Use for complex numbers.
\newcommand{\SYT}{\mathrm{SYT}}
\newcommand{\al}{\alpha}
\newcommand{\cc}{\mathcal{C}}
 \newcommand{\ci}{\mathcal{I}}
 \newcommand{\cj}{\mathcal{J}}
 \newcommand{\cm}{\mathcal{M}}
\newcommand{\cupsigma}{\cc(\gs)}
 \newcommand{\cv}{\mathcal{V}}
\newcommand{\ga}{\alpha}
\newcommand{\gs}{\sigma}
\newcommand{\isigma}{\ci_\gs}
\renewcommand{\ker}{\operatorname{ker }}
\newcommand{\m}[2]{m(#1,#2)}
\newcommand{\mi}[1]{m(#1)}
\newcommand{\pl}{\mathsf{p}}
\newcommand{\rk}{\operatorname{rk}}
\title{Ideals defining components of two-row Springer fibers}
\author{Cristina Sabando-Alvarez, Martha Precup}
\date{Source: arXiv:2606.07507v1 (CC BY 4.0)}
\begin{document}
\maketitle

\noindent\emph{Source and licence.} Ideals defining components of two-row Springer fibers. Version arXiv:2606.07507v1 (CC BY 4.0), distributed by arXiv under the Creative Commons Attribution 4.0 International licence, http://creativecommons.org/licenses/by/4.0/, as recorded in the arXiv licence metadata for this exact version and shown on the abstract page https://arxiv.org/abs/2606.07507v1. Full licence notice: \texttt{licenses/arxiv-2606.07507-CC-BY-4.0.txt}.\par\smallskip

\noindent\textit{Editorial scope.} Counted unit: the article's prop prop.aug.rank (source0.tex lines 2355--2360). The article's complete original proof of it is delivered with it (source0.tex lines 2361--2512, one \texttt{proof} environment of 152 lines). No mathematics is written, repaired or re-derived: the statement and the proof are the article's own lines, in the article's own notation, and no retained line is substituted or re-flowed.  Retained uncounted as the source-local context the proof needs: lines 1190--1194; lines 1327--1332; lines 1574--1578; lines 2039--2047; lines 2106--2107; lines 2311--2312.  Nothing was omitted from any retained span, so the package carries no omission record.  No retained line contains a citation, so the wrapper carries no bibliography and no bibliography entry.  No retained line contains a figure environment, an image-inclusion command or drawing code, so no image is delivered and the package declares no asset dependency.  Editorial formatting only, in addition: an AMS \texttt{amsart} preamble that restates the article's own declarations and macros used by the retained lines, namely the environments \texttt{lemma}, \texttt{definition}, \texttt{prop} on separate counters, together with the standard packages those lines need.  The retained environments print in this document as Lemma 1, Definition 1, Proposition 1 in order of appearance, whereas the article prints the same items with the numbering of its own full document.  The complete native source of the article as distributed in the arXiv e-print for this exact version is bundled unchanged as \texttt{sources/202261/source0.tex}.
\begin{lemma} \label{lem.leftend} For any left end point $i$ of a cup $(i<j)$,
\begin{equation}\label{rem:i}
    i=2S_\gs(i)-\ga(i<j)-1.
\end{equation}
\end{lemma}

\begin{lemma}\label{lem.rankcond} Let $a,k\geq 2$ be a positive integers with $k\geq a+1$ and let $A$ be $2a\times k$ matrix such that 
\begin{equation}\label{vanmins}
\sum_{k=1}^a \pl_{\{k,2a-k+1\},\{c_1,c_2\}}(A)=0 \quad \textup{ for all } \quad 1\leq c_1<c_2\leq k.
\end{equation}
Then $\rk(A)\leq a$, that is, all $(a+1)\times(a+1)$-minors of $A$ are zero.
\end{lemma}

\begin{definition}\label{def:JSigma} Let $\gs$ be a standard tableau of shape $(\ell, \ell)$. We set
\begin{eqnarray*}
\cj_\gs &:=& \sum_{\substack{j\in [n]\\ S_\sigma(j)<\ell}}\left< z_{i,j}, \, z_{i+\ell, j} \mid S_\gs(j)+1\leq i \leq \ell \right> \subset \C[\mathbf{z}].
\end{eqnarray*}
\end{definition}

\begin{lemma}\label{pairingdeterminants} For any positive integer $a\geq 2$ and set of indeterminates $\{v_{i,j}^k \mid i,j\in \{1,2\} \text{ and } 0\leq k \leq a \}$, we have
\begin{eqnarray}\label{lem:ident1}
v_{1,1}^0v_{2,2}^0+\displaystyle\sum_{k=1}^{a-1}\begin{vmatrix} v_{1,1}^k&v_{1,2}^k\\v_{2,1}^k&v_{2,2}^k\end{vmatrix}-v_{2,1}^av_{1,2}^a
=\sum_{k=0}^{a-1}\begin{vmatrix}
v_{1,1}^k&v_{1,2}^{k+1}\\
v_{2,1}^{k+1}&v_{2,2}^k
\end{vmatrix}.
\end{eqnarray}
\end{lemma}

\begin{prop}\label{binv} Let $\gs\in \SYT(\ell, \ell)$. The ideal $\ci_\gs$ is right $B$-invariant. 
\end{prop}

\begin{lemma}\label{lem.kernel} Let $\gs$ be a standard tableau of shape $(\ell,\ell)$ and $\cj_\gs$ the monomial ideal from Definition~\ref{def:JSigma}. We have $A\in \cv(\cj_\gs)$ if and only if $N^{S_\gs(k)} a_k =0$ for all $k$.
\end{lemma}

\begin{prop}\label{prop.aug.rank} Let $\gs$ be a standard tableau of shape $(\ell, \ell)$ and suppose $(i<j)$ is a cup in $\cc(\gs)$ of size at least $2$. If $A\in GL_n(\C)\cap\cv(\ci_\gs)$ then the matrix
\begin{eqnarray}\label{eqn.extmatrix}
\;\;\;\;\;\widetilde{A}_{(i<j)} :=  \begin{bmatrix} | &  & | & | & | &  & |  \\ N^{\mi{i}}a_1 & \cdots & N^{\mi{i}}a_{i-1} & N^{\mi{i}}a_i & N^{\m{i}{i+1}}a_{i+1} & \cdots & N^{\m{i}{j-1}}a_{j-1} \\| &   & | & | & | & &|  \end{bmatrix}
\end{eqnarray}
has $\rk \widetilde{A}_{(i<j)} =\al(i<j)+1$. 
\end{prop}

\begin{proof} Let $A\in GL_n(\C)\cap\cv(\ci_\gs)$ and set $\al=\al(i<j)$. 

It suffices to prove $\rk \widetilde A_{(i<j)} \leq \al+1$. Given this fact, the assertion $\rk \widetilde A_{(i<j)} = \al+1$ follows. Indeed, $N^{\mi{i}}\left( F_i \right) = \C\{ N^{\mi{i}}a_1, \ldots,N^{\mi{i}}a_i \}$ is a subspace of the column space of $\widetilde A_{(i<j)}$ so $\rk \widetilde A_{(i<j)}$ is at least $\dim N^{\mi{i}}\left( F_i \right)$. Since $\dim F_i = i$ and $\dim \ker N^{\mi{i}} = 2\mi{i}$ we have 
\[
\rk \widetilde A_{(i<j)}  \geq \dim N^{\mi{i}}\left( F_i \right) \geq i - 2\mi{i} = \al +1,
\]
where the last equality follows by Lemma~\ref{lem.leftend} using the fact that $\mi{i}=S_\gs(i)-\al-1$. 

We proceed with our proof that $\rk  \widetilde A_{(i<j)} \leq \al+1$. Consider the projection from $\C^n$ to $\C^{2\al+2}$ defined by
\[
\left[ v_1, v_2, \ldots, v_n\right]^T \mapsto \left[v_1,v_2, \ldots, v_{\al+1}, v_{\ell+1}, v_{\ell+2}, \ldots, v_{\ell+\al+1}\right]^T.
\]
By Lemma~\ref{lem.kernel}, we know $a_k\in \ker N^{S_\gs(k)}$ for all $k$. For all $i\leq k \leq j-1$ we have $S_\gs(k) - \m{i}{k} = \al+1$ and thus
\[
N^{\m{i}{k}} a_k \in \ker N^{\al +1} = \C\{e_1, e_2, \ldots, e_{\al+1}\} \oplus \C\{ e_{\ell+1}, e_{\ell+2}, \ldots, e_{\ell+\al+1}\}.
\]
Now consider $k$ with $1\leq k \leq i$. If $S_\gs(k)\leq \mi{i}$ then $N^{\mi{i}}a_k = 0$ and otherwise,
\begin{align}\label{eqn.Nm(i)}
N^{\mi{i}} a_k \in \ker N^{S_\gs(k)-\mi{i}} &= \C\{e_1, e_2, \ldots, e_{S_\gs(k)-\mi{i}}\} \\ \nonumber
& \qquad \qquad \oplus \C\{e_{\ell+1}, e_{\ell+2}, \ldots, e_{\ell + S_\gs(k)-\mi{i}}\}.
\end{align}
In particular, the matrix obtained from $\widetilde A_{(i<j)}$ by replacing each column with its projection to $\C^{2\al+2}$ deletes rows $\al+2, \al+3, \ldots, \ell$ and $\ell+\al+2, \ell+\al+3, \ldots, 2\ell =n$ of $\widetilde A_{(i<j)}$, all of which are zero.  The rank of the resulting $(2\al+2)\times (j-1)$ matrix is equal to the rank of the original. Abusing notation, we denote this matrix by $\widetilde A_{(i<j)}$ as well. 


For any choice of columns $c_1$ and $c_2$ of $\widetilde A_{(i<j)}$ we consider
\begin{eqnarray}\label{eqn.minorsextended}
\pl_{c_1,c_2}\left(\widetilde A_{(i<j)}\right) := \sum_{k=0}^{\al(i<j)}\pl_{\{k+1, 2\al(i<j)+2-k\}, \{c_1,c_2\}} \left(\widetilde A_{(i<j)}\right) .
\end{eqnarray}
The assertion that $\rk \widetilde A_{(i<j)} \leq \al(i<j)+1$ follows directly from Lemma~\ref{lem.rankcond} once we prove~\eqref{eqn.minorsextended} is equal to $0$ for all $1\leq c_1<c_2\leq j-1$. 

If $i\leq c_1<c_2\leq j-1$, then $\pl_{c_1, c_2} = \pl_{c_1, c_2}^{(i<j)}=\sum_{k=0}^{\al}\pl_{\{k+1, 2\al+2-k\}, \{c_1,c_2\}}$ is a generator of $\ci_\gs$, and $\pl_{c_1, c_2}\left(\widetilde A_{(i<j)}\right) = \pl_{c_1, c_2}\left(A\right)=0$ since $A\in \cv(\ci_\gs)$. 

Now suppose $c_1< i <c_2\leq j-1$ and consider $u_{c_1} = I_n + E_{c_1, i} \in B$. Right multiplication of $A$ by $u_{c_1}$ replaces the $i$-th column $a_i$ of $A$ by the sum $a_i+a_{c_1}$. We have 
\[
\pl_{c_1,c_2}\left(\widetilde A_{(i<j)}\right) =  \pl_{i,c_2}^{(i<j)}\left(A\, u_{c_1}\right) -  \pl_{i,c_2}^{(i<j)}\left(A\right)=0-0=0
\]
since $\pl_{i,c_2} = \pl_{i,c_2}^{(i<j)}$ is a generator of $\ci_\gs$, $A\in \cv(\isigma)$, and $Au_{c_1}\in \cv(\ci_\gs)$ as well because $\ci_\gs$ is right $B$-invariant by Proposition~\ref{binv}.


For the rest of the proof, we consider the case of $1\leq c_1<c_2\leq i$.  We start by computing the form of the $c$-th column for some $c\leq i$. First, if $S_\gs(c) \leq \mi{i}$ then $N^{\mi{i}}a_{c} =0$ so the $c$-th column of $\widetilde A_{(i<j)}$ is $0$ and~\eqref{eqn.minorsextended} is obviously $0$ if $c$ is selected as $c_1$ or $c_2$. We may therefore assume that $S_\gs(c) > \mi{i}$. Because $N^{\mi{i}} a_{c}\in \ker N^{S_\gs(c) - \mi{i}}$, the $c$-th column of $\widetilde A_{(i<j)}$ is of the form 
\begin{align} \label{eqn.colform}
\left[ a_{\mi{i}+1, c}\  \right.,  a_{m(i)+2, c}\ , &\cdots ,  a_{\mi{i}+(S_\gs(c)-\mi{i}), c}\ , 0^{S_\gs(i) - S_\gs(c)} \ ,   \\  
\nonumber &  a_{\ell+\mi{i}+1, c} \ ,    a_{\ell + m(i)+2, c}\ , \cdots , a_{\ell + \mi{i}+(S_\gs(c)-\mi{i}), c}\ ,\left. 0^{S_\gs(i) - S_\gs(c)}  \right]^T.
\end{align}
Note that this vector has $2\al(i<j)+2$ many entries since $\mi{i} = S_\gs(i)-\alpha(i<j)-1$ and thus 
\[
(S_\gs(c)-\mi{i}) + (S_\gs(i) -  S_{\gs}(c) )  =  \alpha(i<j)+1.
\]
Now, for $c_1<c_2\leq i$, the equation~\eqref{eqn.minorsextended} becomes
\begin{eqnarray}\label{eqn.minorsextended2}
\pl_{c_1,c_2}\left(\widetilde A_{(i<j)}\right) = \sum_{k=0}^{\al(i<j)} \begin{vmatrix} a_{\mi{i}+k+1, c_1} & a_{\mi{i}+k+1, c_2}\\ a_{\ell+\mi{i}+\al(i<j) +1 - k, c_1} & a_{\ell+\mi{i}+\al(i<j) +1 - k, c_2}  \end{vmatrix} .
\end{eqnarray}


Suppose $S_\gs(c_1)+S_\gs(c_2)\leq i$. Under this condition, we claim that every minor in the sum~\eqref{eqn.minorsextended2} is $0$. Indeed, since the $c$-th column has form~\eqref{eqn.colform}, for all $0\leq k \leq \al(i<j)$ we have 
\begin{align}\label{eqn.ineq1}
a_{\mi{i}+k+1, c} \neq 0 & \Rightarrow  \mi{i}+1 \leq \mi{i} +k + 1 \leq \mi{i}+(S_\gs(c)-\mi{i}) \\
\nonumber &\Rightarrow k \leq S_\gs(c)-\mi{i}  -1
\end{align}
and 
\begin{align} \label{eqn.ineq2}
&a_{\ell+\mi{i}+\al(i<j) +1 - k, c} \neq 0 \\
\nonumber  & \qquad \qquad \Rightarrow  \mi{i}+1  \leq  \mi{i} + \alpha(i<j)+1 - k \leq  \mi{i} + (S_\gs(c)-\mi{i}) \\
\nonumber & \qquad \qquad \Rightarrow  \mi{i}+ \al(i<j) +1 - S_\gs(c) \leq k.
\end{align}
Thus, using~\eqref{eqn.ineq1} for $c=c_1$ and~\eqref{eqn.ineq2} for $c=c_2$ we have
\begin{align*}
a_{\mi{i}+k+1, c_1}&a_{\ell+\mi{i}+\al(i<j) +1 - k, c_2} \neq 0 \\ &\Rightarrow   \mi{i}+ \al(i<j) +1 - S_\gs(c_2)  \leq S_\gs(c_1)-\mi{i}  -1 \\
&\Rightarrow  2\mi{i} +\al (i<j) +1 \leq S_\gs(c_1)+S_\gs(c_2) -1 \\
&\Rightarrow i  < S_\gs(c_1)+S_\gs(c_2),
\end{align*}
where the last equality follows since $2\mi{i} +\al (i<j) +1 = 2S_\gs(i) -\al(i<j)-1 = i$ by Lemma~\ref{lem.leftend}. By the same reasoning with the roles of $c_1$ and $c_2$ reversed we have 
\[
a_{\mi{i}+k+1, c_2}a_{\ell+\mi{i}+\al(i<j) +1 - k, c_1} \neq 0 \Rightarrow i  < S_\gs(c_1)+S_\gs(c_2).
\]
Thus, if $S_\gs(c_1)+S_\gs(c_2)\leq i$, we have shown that  
\[
\begin{vmatrix} a_{\mi{i}+k+1, c_1} & a_{\mi{i}+k+1, c_2}\\ a_{\ell+\mi{i}+\al(i<j) +1 - k, c_1} & a_{\ell+\mi{i}+\al(i<j) +1 - k, c_2}  \end{vmatrix} =0
\]
for all $0\leq k \leq \al(i<j)$. This proves~\eqref{eqn.minorsextended2} is $0$ in this case since every minor in the sum is $0$, proving our claim. 






Now assume $S_\gs(c_1) + S_\gs(c_2)>i$.  Using Lemma~\ref{lem.leftend} once more, we have 
\begin{align*}
i<S_\gs(c_1) + S_\gs(c_2) &\Leftrightarrow i + 1 \leq  S_\gs(c_1) + S_\gs(c_2) \\
& \Leftrightarrow 2 S_{\gs}(i)-\al(i<j) \leq S_\gs(c_1) + S_\gs(c_2) \\
& \Leftrightarrow 2 S_{\gs}(i)- S_\gs(c_1)-S_\gs(c_2) \leq \al(i<j). 
\end{align*}
Since $c_1<c_2\leq i$ we know $S_\gs(c_1)\leq S_\gs(c_2)\leq S_\gs(i)$ and thus $2 S_{\gs}(i)- S_\gs(c_1)-S_\gs(c_2)\geq 0$. For the remainder of the proof, let $(i'<j')$ be the unique cup above $(i<j)$ such that 
\begin{eqnarray}\label{eqn.i'<j'def}
\al(i'<j') = \al(i<j) - \left( 2 S_{\gs}(i)- S_\gs(c_1)-S_\gs(c_2) \right) \geq 0.
\end{eqnarray}
Since the cup $(i'<j')$ has $2 S_{\gs}(i)- S_\gs(c_1)-S_\gs(c_2)$ fewer cups nested above it than $(i<j)$, 
\begin{align*}
S_\gs(i') \leq  S_\gs(i) - \left(2 S_{\gs}(i)- S_\gs(c_1)-S_\gs(c_2) \right) \Rightarrow S_\gs(i')+S_\gs(i) \leq S_\gs(c_1)+S_\gs(c_2).
\end{align*}
Because $S_\gs(i)\geq S_\gs(c_2)$, the last inequality implies $S_\gs(i') \leq S_\gs(c_1)$. Furthermore, as $i'$ is the left endpoint of a cup, it is the minimal index with cup sequence value $S_\gs(i')$, and we conclude $i'\leq c_1$. Thus $c_1$ and $c_2$ are vertices of the cup diagram $\cupsigma$ nested inside the cup $(i'<j')$, and they index columns in the matrix $A_{(i'<j')}$. 


Since $S_\gs(c)-\mi{i}-1 = \al(i<j) - \left( S_\gs(i) - S_\gs(c_1) \right)$ and $\mi{i}+\al(i<j) +1 -S_\gs(c_1) = S_\gs(i)-S_{\gs}(c_2)$, using  equations~\eqref{eqn.ineq1} and~\eqref{eqn.ineq2} for $c_1$ and $c_2$ yields the inequalities:
\begin{align}
\label{eqn.ineqc1A} a_{\mi{i}+k+1, c_1} \neq 0 &\Rightarrow k \leq   \al(i<j) - \left(S_\gs(i)-S_\gs(c_1)\right) \\
\label{eqn.ineqc2A} a_{\mi{i}+k+1, c_2} \neq 0 &\Rightarrow k \leq   \al(i<j) - \left(S_\gs(i)-S_\gs(c_2)\right) \\
\label{eqn.ineqc1B} a_{\ell + \mi{i} + \al(i<j)+1-k,c_1} \neq 0 &\Rightarrow S_\gs(i)-S_\gs(c_1) \leq k \\
\label{eqn.ineqc2B} a_{\ell + \mi{i} + \al(i<j)+1-k,c_2} \neq 0 &\Rightarrow S_\gs(i)-S_\gs(c_2) \leq k .
\end{align}
Note that $c_1< c_2$ implies $S_\gs(c_1)\leq S_\gs(c_2)$. By~\eqref{eqn.ineqc1A} and ~\eqref{eqn.ineqc2A}, we have 
\[
a_{\mi{i}+k+1, c_1} = a_{\mi{i}+k+1, c_2} = 0 \; \text{ when } \; k > \al(i<j) - \left( S_\gs(i)-S_{\gs(c_2)} \right)
\]
and by~\eqref{eqn.ineqc1B} and ~\eqref{eqn.ineqc2B}, we have  
\[
a_{\ell + \mi{i} + \al(i<j)+1-k,c_1} = a_{\ell + \mi{i} + \al(i<j)+1-k,c_2} = 0 \; \text{ when } \; k< S_{\gs(i)} - S_\gs(c_2).
\]
Thus, the only minors in the sum~\eqref{eqn.minorsextended2} which are nonzero occur for $k$ such that $S_\gs(i)-S_\gs(c_2) \leq k \leq \al(i<j)- \left( S_\gs(i)-S_\gs(c_2) \right)$. Equation~\eqref{eqn.minorsextended2} becomes
\begin{align}\label{eqn.minorsextended3}
\pl_{c_1, c_2}(\widetilde A_{(i<j)}) & =  \sum_{k=S_\gs(i)-S_\gs(c_2)}^{\al(i<j) - \left( S_\gs(i)-S_\gs(c_2) \right)} \begin{vmatrix} a_{\mi{i}+k+1, c_1} & a_{\mi{i}+k+1, c_2}\\ a_{\ell+\mi{i}+\al(i<j) +1 - k, c_1} & a_{\ell+\mi{i}+\al(i<j) +1 - k, c_2}  \end{vmatrix}. 
\end{align}

Let $t(c_1, c_2):= S_\gs(c_2)-S_{\gs}(c_1) \geq 0$. 
By~\eqref{eqn.ineqc1B}, $a_{\ell + \mi{i} + \al(i<j)+1-k,c_1}=0$ for all $k< S_\gs(i)-S_\gs(c_1)$, i.e., for all $k$ such that $S_\gs(i)-S_\gs(c_2)\leq k \leq S_\gs(i)-S_\gs(c_2) +t(c_1,c_2)-1$. Similarly, by~\eqref{eqn.ineqc1A}, $a_{\mi{i}+k+1, c_1}=0$ for all $k> \al(i<j) - \left( S_\gs(i) - S_\gs(c_1)\right)$, i.e., for all $k$ such that 
\begin{align*}
\al(i<j) - \left( S_\gs(i) - S_\gs(c_1)\right) + 1\leq k &\leq \al(i<j) - \left( S_\gs(i) - S_\gs(c_1)\right) +t(c_1, c_2)\\ & \qquad \qquad \qquad  =  \al(i<j) - \left( S_\gs(i) - S_\gs(c_2)\right).
\end{align*}
This means we can further decompose~\eqref{eqn.minorsextended3} as: 
\begin{align*}
\pl_{c_1, c_2}(\widetilde A_{(i<j)}) & =   \left( \sum_{k=S_\gs(i)-S_\gs(c_2)}^{S_\gs(i) -S_\gs(c_2) + t(c_1,c_2)-1}  a_{\mi{i}+k+1, c_1}a_{\ell+\mi{i}+\al(i<j) +1 - k, c_2} \right)  \\
& \quad \quad + \left( \ \sum_{k=S_\gs(i)-S_\gs(c_2) + t(c_1, c_2)}^{\al(i<j)-\left( S_\gs(i)-S_\gs(c_1) \right)} \begin{vmatrix} a_{\mi{i}+k+1, c_1} & a_{\mi{i}+k+1, c_2}\\ a_{\ell+\mi{i}+\al(i<j) +1 - k, c_1} & a_{\ell+\mi{i}+\al(i<j) +1 - k, c_2}  \end{vmatrix} \ \right) \\
& \qquad \quad \quad  -  \left( \sum_{k=\al(i<j) - \left( S_\gs(i)-S_\gs(c_1) \right)+1 }^{\al(i<j) - \left( S_\gs(i)-S_\gs(c_2) \right)} a_{\mi{i}+k+1, c_2}a_{\ell+\mi{i}+\al(i<j) +1 - k, c_1}  \right).
\end{align*}
If $t(c_1,c_2) = 0$, then the first and last expressions above are both $0$, otherwise there are $t(c_1,c_2)$ terms in each of the first and last sums. 
We simplify this expression by applying Lemma~\ref{pairingdeterminants} exactly $t(c_1, c_2)$ many times to obtain:
\begin{align*}\label{eqn.minorsextended4}
\pl_{c_1, c_2}(\widetilde A_{(i<j)}) & =   \sum_{k=S_\gs(i)-S_\gs(c_2)}^{\al(i<j)-\left( S_\gs(i)-S_\gs(c_1) \right)} \begin{vmatrix} a_{\mi{i}+k+1, c_1} & a_{\mi{i}+k+t(c_1,c_2)+1, c_2}\\ a_{\ell+\mi{i}+\al(i<j) +1 - \left( k + t(c_1,c_2) \right) , c_1} & a_{\ell+\mi{i}+\al(i<j) +1 - k, c_2}  \end{vmatrix}.
\end{align*}
Finally, we re-index the sum over $0\leq k \leq \al(i<j) - \left( 2S_\gs(i) - S_\gs(c_1)-S_\gs(c_2) \right) = \al (i'<j')$. Using the definition of $\al(i'<j')$ from~\eqref{eqn.i'<j'def} and the fact that $\mi{i} = S_\gs(i)-\al(i<j)-1$ we get  
\begin{align*}
S_\gs(c_1) - \al(i'<j')+k & = \mi{i} + k + S_\gs(i)-S_{\gs}(c_2)+1  \\
S_\gs(c_2) - \al(i'<j')+k &= \mi{i} + k + S_\gs(i)-S_{\gs}(c_2) + t(c_1, c_2)+1 \\
S_\gs(c_2)+k & = \mi{i} +\al(i<j) +1 - \left( k + S_\gs(i)-S_\gs(c_2) \right) \\
S_\gs(c_1)+k & = \mi{i} +\al(i<j) +1 - \left( k + S_\gs(i)-S_\gs(c_2) +t(c_1,c_2)  \right) \ ,
\end{align*}
and reindexing gives us 
\[
\pl_{c_1, c_2}(\widetilde A_{(i<j)})  = \sum_{k=0}^{\al(i'<j')} \begin{vmatrix} a_{S_\gs(c_1)-\al(i'<j')+k, c_1} &  a_{S_\gs(c_2)-\al(i'<j')+k, c_2} \\  a_{\ell-S_\gs(c_1)-k,c_1}  & a_{\ell + S_\gs(c_2) - k, c_2 }\end{vmatrix} = \pl_{c_1, c_2}^{(i'<j')} (A).
\]
Thus, $\pl_{c_1, c_2}(\widetilde A_{(i<j)}) = \pl_{c_1, c_2}^{(i'<j')} (A)=0$ since $\pl_{c_1, c_2}^{(i'<j')}$ is a generator of $\cm_{\gs, (i'<j')} \subseteq \ci_\gs$ and $A\in \cv(\ci_\gs)$. This completes the proof. 
\end{proof}

\end{document}
